\chapter{Vivado学习记录}
\section{Vivado 使用操作}
\subsection{工程操作}
\begin{enumerate}
    \item 打开已有工程：Open Project打开.xpr文件
    \item 清理vivado工程：
\end{enumerate}

\subsection{调试操作}
\begin{enumerate}
    \item 出现syntax error表示文件出现语法错误
\end{enumerate}

\section{代码记录}

\begin{minted}{vhdl}
--这是单行注释
/*
这是
多行注释
VHDL基本语法结构
*/

--VHDL不区分大小写
--library 
--library的声明
LIBRARY library_name;
USE ieee.std_logic_1164.all;--std_logic_1164是package_name,all是package_parts
--entity
ENTITY test_logic is
    PORT(
        a,b,c,d :IN STD_LOGIC;--STD_LOGIC是IN的数据类型
        and_out,or_out;OUT STD_LOGIC --PORT里面最后一行不用加分号;
    );
end entity test_logic;

--architecture
architecture bhv of test_logic is
--[declarations]这里的声明部分为可选，用于对信号和常量等进行声明
BEGIN
    and_out<= a AND b;
    or_out <=c Or d;
end architecture bhv ;
\end{minted}

\begin{minted}{verilog}
/*
verilog中是模块体
assign语句是赋值语句，对wire类型数据有用
*/
module test(CLK,D,Q);//模块名后紧跟端口名
    output Q;
    input CLK ,D;
    reg Q;//寄存器类型
    always @(event)//敏感列表，()里只要有变化就执行begin的内容
/*包括边沿敏感和电平敏感,posedge上升沿到来,negedge下降沿到来
always只对reg的数据类型起作用*/
    begin
        Q<=D;
    end

endmodule
\end{minted}

\begin{myminted}[title=状态机语法]{verilog}
/*完整的状态机包含：
状态寄存器
次态逻辑
输出逻辑*/
parameter =;
\end{myminted}